% Chapter 6, typeset from lectures/L06: the README is the opener and the closing
% review, and each appendix is a section, in the same order and with the same subsections.

\chapter{Subsystems and Debugging}
\label{c6:ch:subsystems}\label{c6:ch}

{\large\itshape Making a program talk back, remember things, and crash usefully.\par}

\begin{chapterbox}{In this chapter}
\begin{itemize}
  \item The shell: built-in commands, commands of your own with subcommands and arguments, and changing log
    levels at run time.
  \item The settings subsystem: key-value pairs that survive a restart, stored with ZMS on flash.
  \item \code{native_sim}’s simulated flash: a file, and the command-line options that control it.
  \item Networking on \code{native_sim} through offloaded sockets: Zephyr’s socket API on the host’s
    sockets.
  \item Fatal errors: \code{k_oops()}, \code{k_panic()}, assertions, and \code{k_sys_fatal_error_handler()}.
  \item The host’s tools on a Zephyr program: AddressSanitizer and UndefinedBehaviorSanitizer.
  \item Debugging a \code{native_sim} program with GDB.
  \item Code coverage with gcov and Twister.
  \item What moving to real hardware would change, and what it would not.
\end{itemize}
\end{chapterbox}

This last chapter is about the parts of Zephyr that make a program usable by someone other than its
author. A shell, so a running system can be asked what it is doing. Settings that survive a restart,
on the flash that \code{native_sim} keeps in a file. Sockets, which on this board are the host's
own, so a Zephyr program can be talked to from an ordinary terminal.

Then it is about what to do when the program is wrong. Zephyr's fatal error handling decides what an
impossible value costs, and the host's tools, AddressSanitizer, UndefinedBehaviorSanitizer, GDB and
gcov, work on a \code{native_sim} program exactly as on any other Linux program. That is the best
argument this book has for keeping code runnable on \code{native_sim} even when the product is a
board.

% Section 6.1, from lectures/L06/appendix/a_subsystems.md.

\section{Subsystems: the Shell, Settings and Networking}
\label{c6:sec:subsystemstheshell}\label{c6:app:a}
Chapter~\ref{c1:ch} to Chapter~\ref{c5:ch} built programs that print. This section covers three
subsystems that make a program part of a system: a shell to talk to it, settings to make it
remember, and sockets to put it on a network. All three are the same code on a board and on
\code{native_sim}; what \code{native_sim} puts under them, a terminal, a file and the host’s
network, is what this section is careful to separate from the subsystems themselves. By the end you
should be able to explain every file in the \file{lectures/L06/examples/shell},
\file{lectures/L06/examples/settings} and \file{lectures/L06/examples/sockets} examples.

\subsection{The shell}

\subsubsection{What it is}
Zephyr’s shell is a command line inside the firmware: line editing, history, \code{Tab} completion,
a \code{help} for every command, and colors. Subsystems register commands of their own, so enabling
the shell alone already gives you commands for the kernel, and enabling logging or settings with it
gives you commands for those. An application adds its own the same way.

\begin{kconfigcode}
CONFIG_SHELL=y
\end{kconfigcode}

The shell runs in a thread of its own, and reads and writes through a \textbf{backend}. The usual
one is the serial backend, on the UART that the devicetree’s \code{chosen} node names
\code{zephyr,shell-uart}; \code{native_sim}’s devicetree names \code{uart0}. Others carry the same
shell over SEGGER RTT, telnet, a WebSocket, MQTT or RPMsg, and a \emph{dummy} backend runs commands
from code, which is how the shell is tested. The thread of the serial backend is called
\code{shell_uart}.

\subsubsection{The shell on \code{native_sim}}
\code{native_sim}’s \code{uart0} is a pseudo-terminal by default. The program prints \code{uart
connected to pseudotty: /dev/pts/11} (the number varies), and you connect to that from a second
terminal. For a course, and for tests, it is simpler to put the UART on the terminal the program was
started from:

\begin{kconfigcode}
CONFIG_UART_NATIVE_PTY_0_ON_STDINOUT=y
\end{kconfigcode}

or, without rebuilding, the command-line option \code{-uart_stdinout}. Commands can then be typed,
or piped in, which is how a session is scripted:

\begin{shell}
printf 'greet Ada\nled on\nled status\n' | ./build/zephyr/zephyr.exe -stop_at=1
\end{shell}

The shell echoes each command after its prompt, \code{uart:~$}, and prints the response on the lines
that follow.

\subsubsection{Commands}
A command is a handler function and a static entry that describes it. Every handler has the same
signature:

\begin{ccode}
static int cmd_greet(const struct shell *sh, size_t argc, char **argv)
{
	shell_print(sh, "Hello, %s!", argv[1]);
	return 0;
}
\end{ccode}

\begin{itemize}
  \item \textbf{\code{sh}} is the shell instance that ran the command, and the one its output must go to:
    there may be several backends, each with a shell of its own.
  \item \textbf{\code{argv}} is the command line split into words, \textbf{\code{argv[0]}} being the
    command’s own name, and \textbf{\code{argc}} is how many there are.
  \item \textbf{The return value} is \code{0} on success and a negative \code{errno} on failure.
\end{itemize}

A command with no subcommands is registered with its argument counts:

\begin{ccode}
SHELL_CMD_ARG_REGISTER(greet, NULL, "Greet someone: greet <name> [times]", cmd_greet, 2, 1);
\end{ccode}

The last two numbers are the \textbf{mandatory} and \textbf{optional} word counts, and the mandatory
count \textbf{includes the command’s own name}: \code{2, 1} means \code{greet} and a name, and
optionally one more word. The shell checks the count before it calls the handler, so the handler
never sees \code{argc} outside the range. A wrong count prints the command’s help instead:

\begin{console}
uart:~$ greet
greet: wrong parameter count
greet - Greet someone: greet <name> [times]
\end{console}

\code{SHELL_CMD_REGISTER()} is the same without counts, for a command that takes no arguments or
checks them itself. Commands with \textbf{subcommands} have a table of them, built at compile time:

\begin{ccode}
SHELL_STATIC_SUBCMD_SET_CREATE(sub_led,
	SHELL_CMD(on, NULL, "Turn the LED on.", cmd_led_on),
	SHELL_CMD(off, NULL, "Turn the LED off.", cmd_led_off),
	SHELL_CMD_ARG(blink, NULL, "Blink: blink <times>", cmd_led_blink, 2, 0),
	SHELL_SUBCMD_SET_END);

SHELL_CMD_REGISTER(led, &sub_led, "Control the LED.", NULL);
\end{ccode}

A subcommand is written with \code{SHELL_CMD()} or, with counts, \code{SHELL_CMD_ARG()}, and its
counts again include its own name: in \code{led blink 3}, the handler for \code{blink} gets
\code{argv[0] = "blink"} and \code{argv[1] = "3"}. A command with subcommands and no handler of its
own, \code{NULL} above, prints its help and the list of subcommands when it is run alone. Because
the tables are static, \code{Tab} can complete them. (\code{SHELL_DYNAMIC_CMD_CREATE()} creates
subcommands at run time, such as a list of device names, and \code{SHELL_COND_CMD_REGISTER()}
registers a command only when a Kconfig option is enabled.)

\subsubsection{Output}
A handler prints through the shell it was called by, never with \code{printk()}, whose output would
go to the console and not necessarily to the backend the user is on:

\begin{booktable}{@{}LL@{}}
  \thead{Function} & \thead{Output} \\
  \midrule
  \code{shell_print(sh, fmt, ...)} & Normal output. \\
  \code{shell_info(sh, fmt, ...)} & Informational output, in green. \\
  \code{shell_warn(sh, fmt, ...)} & A warning, in yellow. \\
  \code{shell_error(sh, fmt, ...)} & An error, in red. \\
  \code{shell_hexdump(sh, data, len)} & A buffer, as a hex dump. \\
\end{booktable}

The functions add the newline themselves.

\subsubsection{Validating arguments}
The shell counts words; what the words say is the handler’s business. \code{atoi()} is the wrong
tool, because it cannot fail: \code{atoi("warm")} is 0 and \code{atoi("25x")} is 25. \code{strtol()}
reports where it stopped and whether the value overflowed, so a word is a number exactly when it
stopped at the end of the word, after at least one digit, with \code{errno} still 0:

\begin{ccode}
char *end;

errno = 0;
long value = strtol(argv[1], &end, 10);

if ((end == argv[1]) || (*end != '\0') || (errno != 0)) {
	shell_error(sh, "not a temperature: '%s'", argv[1]);
	return -EINVAL;
}
\end{ccode}

A value that is a number but out of range is a different error, conventionally \code{-ERANGE}.

\subsubsection{Return values and \code{retval}}
The shell does not print a handler’s return value, but it keeps it: the built-in command
\code{retval} prints the value the previous command returned. That makes errors visible to a script,
and makes the shell’s own conventions visible to you:

\begin{console}
uart:~$ greet Ada 9
times must be between 1 and 5, not '9'
uart:~$ retval
-22
uart:~$ foo
foo: command not found
uart:~$ retval
-8
\end{console}

\code{-22} is \code{-EINVAL}, returned by the handler; \code{-8} is \code{-ENOEXEC}, returned by the
shell for a command that does not exist. A command with subcommands run without one returns \code{1}
(\code{SHELL_CMD_HELP_PRINTED}), since all it did was print its help.

\subsubsection{The built-in commands}
With the shell come commands from the subsystems that are enabled. In Zephyr 4.4:

\begin{booktable}{@{}LLL@{}}
  \thead{Command} & \thead{From} & \thead{What it does} \\
  \midrule
  \code{help}, \code{history}, \code{clear}, \code{retval} & The shell itself, \code{help} always
    and the others by default & The shell’s own commands. \\
  \code{kernel uptime}, \code{kernel version} & The kernel shell (\code{CONFIG_KERNEL_SHELL}, on
    with the shell) & The system’s uptime and Zephyr’s version. \\
  \code{kernel thread list}, \code{kernel thread stacks} & The kernel shell & Every thread, its
    priority, state and stack usage. Zephyr 3.7 called them \code{kernel threads} and
    \code{kernel stacks}. \\
  \code{log status}, \code{log enable <level> <module>}, \code{log disable <module>} & Logging, with
    \code{CONFIG_LOG_CMDS=y} & Every log module’s level, changed at run time. \\
  \code{settings list}, \code{settings read <name>}, \code{settings write}, \code{settings delete} &
    Settings, with \code{CONFIG_SETTINGS_SHELL=y} & The stored settings. \\
\end{booktable}

Two of them deserve a warning on \code{native_sim}:

\begin{itemize}
  \item \textbf{\code{kernel thread list}} prints stack usage such as \code{usage 40 / 3072 (1 \%)}, which
    means nothing here: Chapter~\ref{c3:ch} showed that threads run on host stacks. On a board it is how
    stacks are sized.
  \item \textbf{\code{log enable dbg app}} answers \code{app: level set to inf.} for a module compiled at
    \code{LOG_LEVEL_INF}. A level can be lowered at run time, never raised above what was compiled in,
    because the messages above it are not in the program at all (Chapter~\ref{c1:ch}).
\end{itemize}

\subsection{Settings}

\subsubsection{What it is}
The settings subsystem stores \textbf{key-value pairs} in non-volatile memory, so that a device
remembers them across a restart: a calibration, a network name, a setpoint. The keys are names made
of parts separated by \code{/}, such as \code{thermostat/setpoint}, and the first part is a
\textbf{subtree} that one module owns. The values are bytes, of whatever size and type the owner
decides.

\begin{kconfigcode}
CONFIG_SETTINGS=y
\end{kconfigcode}

\noindent The subsystem does not write to memory itself. It hands the pairs to a \textbf{backend},
chosen with a Kconfig \code{choice}. Its defaults include, in order:

\begin{booktable}{@{}LL@{}}
  \thead{Backend} & \thead{Chosen by default when} \\
  \midrule
  ZMS, Zephyr Memory Storage & \code{CONFIG_ZMS=y} \\
  NVS, Non-Volatile Storage & \code{CONFIG_NVS=y} \\
  FCB, Flash Circular Buffer & \code{CONFIG_FCB=y} \\
  A file & \code{CONFIG_FILE_SYSTEM=y} \\
  None & nothing else is enabled: settings then live in RAM only \\
\end{booktable}

ZMS and NVS are both log-structured key-value stores for flash: they append new values and erase
whole sectors when the old values are no longer needed, so that no single flash cell wears out
early. ZMS is the newer of the two, and is designed for memories that need no erase as well. The
course uses ZMS:

\begin{kconfigcode}
CONFIG_FLASH=y
CONFIG_FLASH_MAP=y
CONFIG_ZMS=y
CONFIG_SETTINGS=y
\end{kconfigcode}

\noindent \textbf{Where on the flash} is a devicetree question. The ZMS backend uses the partition
that the \code{chosen} node names \code{zephyr,settings-partition}, or, when there is none, the
partition labelled \code{storage_partition}. \code{native_sim}’s devicetree has one: 16 KiB at
offset \code{0xfc000} of its 2 MiB flash.

\subsubsection{Handlers}
A module that owns a subtree registers a \textbf{handler} for it: a set of callbacks, any of which
may be \code{NULL}.

\begin{booktable}{@{}LLL@{}}
  \thead{Callback} & \thead{Called by} & \thead{To} \\
  \midrule
  \code{h_set(name, len, read_cb, cb_arg)} & \code{settings_load()}, once for every stored setting
    in the subtree & Take the stored value into the module’s variables. \\
  \code{h_commit()} & \code{settings_load()}, after every \code{h_set} & Act on the whole set of
    values at once. \\
  \code{h_export(export_func)} & \code{settings_save()} & Hand every value of the subtree over for
    saving. \\
  \code{h_get(name, val, len)} & \code{settings_runtime_get()} & Return the value the module holds
    in RAM. \\
\end{booktable}

The handler is defined at compile time, like a shell command:

\begin{ccode}
SETTINGS_STATIC_HANDLER_DEFINE(thermostat, "thermostat", NULL, thermostat_set, NULL, NULL);
\end{ccode}

\noindent \code{h_set} gets the rest of the name, after the subtree, and \textbf{not the value
itself}. The subsystem does not know how big the value should be, so the handler checks the stored
length and reads the value into storage of its own:

\begin{ccode}
static int thermostat_set(const char *name, size_t len, settings_read_cb read_cb, void *cb_arg)
{
	const char *next;

	if (!settings_name_steq(name, "setpoint", &next) || (next != NULL)) {
		return -ENOENT;
	}
	if (len != sizeof(setpoint_c)) {
		return -EINVAL;
	}

	ssize_t read = read_cb(cb_arg, &setpoint_c, sizeof(setpoint_c));

	return (read < 0) ? (int)read : 0;
}
\end{ccode}

\noindent \code{settings_name_steq(name, "setpoint", &next)} is true when \code{name} begins with
\code{setpoint}, and sets \code{next} to what follows the next \code{/}, or to \code{NULL} if
nothing does. The length check is not pedantry: flash written by an older firmware, with a smaller
type, is still there after an update.

\subsubsection{Loading, saving and deleting}
\begin{ccode}
settings_subsys_init();                       /* Mount the backend. */
settings_load();                              /* Call h_set for everything stored, then h_commit. */
settings_save_one("thermostat/setpoint", &value, sizeof(value));
settings_delete("thermostat/setpoint");
\end{ccode}

\textbf{Loading happens once, at boot, and on the application’s initiative.} Until
\code{settings_load()} runs, every variable holds its initial value, which is also what it keeps if
nothing was ever stored: on the very first boot, no handler is called at all. Saving writes one
value at once; there is no transaction around several.

\subsubsection{The flash on \code{native_sim}}
\code{native_sim}’s flash is simulated in a file, \textbf{\code{flash.bin}, in the directory the
program runs in}:

\begin{shell}
west build -t run          # Runs in build/, so the file is build/flash.bin.
./build/zephyr/zephyr.exe  # Runs where you are, and uses a flash.bin there.
\end{shell}

\noindent The file is the whole 2 MiB of the simulated flash. A pristine build, \code{west build
-p}, deletes \code{build/} and the flash with it: a factory-new device. The command-line options
decide what happens to the flash for one run:

\begin{booktable}{@{}lL@{}}
  \thead{Option} & \thead{Effect} \\
  \midrule
  \code{-flash=<path>} & Use another file. \\
  \code{-flash_erase} & Erase the flash at startup. \\
  \code{-flash_rm} & Delete the file when the program ends. \\
  \code{-flash_in_ram} & Keep the flash in RAM only. Every run starts erased, and the file is left
    alone. \\
\end{booktable}

\code{strings build/flash.bin} shows the names of the stored settings, \code{boot/count} for the
\file{lectures/L06/examples/settings} example, because the backend stores each name next to its
value. With \code{CONFIG_SETTINGS_SHELL=y}, the shell shows the same from inside:

\begin{console}
uart:~$ settings list
thermostat/setpoint
uart:~$ settings read thermostat/setpoint
00000000: 1b 00 00 00                                      |....             |
\end{console}

\noindent \code{1b 00 00 00} is 27, as a little-endian \code{int32_t}: the bytes exactly as the
handler wrote them.

\subsection{Networking on \code{native_sim}}

\subsubsection{Two ways onto a network}
Zephyr has its own TCP/IP stack, and on a board that is what runs under the socket API.
\code{native_sim} can run it too, over a TAP interface on the host (\code{CONFIG_ETH_NATIVE_TAP}),
but that needs the TAP interface set up on the host, with root privileges.

The alternative is \textbf{offloaded sockets}:

\begin{kconfigcode}
CONFIG_NETWORKING=y
CONFIG_NET_IPV4=y
CONFIG_NET_SOCKETS=y
CONFIG_NET_DRIVERS=y
CONFIG_NET_SOCKETS_OFFLOAD=y
CONFIG_NET_NATIVE_OFFLOADED_SOCKETS=y
\end{kconfigcode}

\noindent With them, every socket call is handed to the host’s own socket API. A server listens on a
real port of your machine, a client connects to real addresses, and nothing needs setting up. The
price is in the name: Zephyr’s TCP/IP stack is not in the program at all, and Linux does the
networking. That is what you want while developing the protocol on top, and it is exactly why the
same program still has to be tested on a board.

\subsubsection{The socket API}
Zephyr’s socket API is BSD sockets with a prefix:

\begin{booktable}{@{}LL@{}}
  \thead{BSD} & \thead{Zephyr} \\
  \midrule
  \code{socket()}, \code{bind()}, \code{listen()}, \code{accept()}, \code{connect()} &
    \code{zsock_socket()}, \code{zsock_bind()}, \code{zsock_listen()}, \code{zsock_accept()},
    \code{zsock_connect()} \\
  \code{send()}, \code{recv()}, \code{sendto()}, \code{recvfrom()}, \code{close()} &
    \code{zsock_send()}, \code{zsock_recv()}, \code{zsock_sendto()}, \code{zsock_recvfrom()},
    \code{zsock_close()} \\
  \code{setsockopt(s, SOL_SOCKET, SO_RCVTIMEO, ...)} &
    \code{zsock_setsockopt(s, ZSOCK_SOL_SOCKET, ZSOCK_SO_RCVTIMEO, ...)} \\
  \code{struct sockaddr_in}, \code{AF_INET}, \code{htons()} & \code{struct net_sockaddr_in},
    \code{NET_AF_INET}, \code{net_htons()} \\
\end{booktable}

\textbf{Zephyr 4.4 prefixed the types and constants too.} Before it, \code{struct sockaddr_in} and
\code{AF_INET} were Zephyr’s own names, which collided with the C library’s and POSIX’s. Now the
networking headers define \code{net_}, \code{NET_} and \code{ZSOCK_} names, and a compatibility
header, enabled by \code{CONFIG_NET_NAMESPACE_COMPAT_MODE} (on by default), keeps the old names
working. New code should use the prefixed ones, as the examples do; code that wants the plain POSIX
names includes the POSIX headers and enables Zephyr’s POSIX API.

\subsubsection{Blocking}
A blocking call blocks \textbf{the calling thread}, not the program: \code{zsock_recv()} waits
inside the kernel, like \code{k_sem_take()}, and every other thread runs meanwhile. Underneath, the
offloaded sockets driver asks the host which sockets are ready once every simulated millisecond
(\code{CONFIG_NET_NATIVE_OFFLOADED_SOCKETS_EPOLL_WAIT_INTERVAL}).

That holds only while the driver knows the socket’s timeouts. \textbf{In Zephyr 4.4.2, the driver
allocates the socket that \code{zsock_accept()} returns without initializing its timeouts}, and a
blocking \code{zsock_recv()} on it can go straight to the host’s \code{recv()}, which blocks the
whole process, simulated time and every thread with it. The \file{lectures/L06/examples/sockets}
example froze at \code{Client connected} until it set both timeouts explicitly after accepting:

\begin{ccode}
struct zsock_timeval forever = {0};

zsock_setsockopt(client, ZSOCK_SOL_SOCKET, ZSOCK_SO_RCVTIMEO, &forever, sizeof(forever));
zsock_setsockopt(client, ZSOCK_SOL_SOCKET, ZSOCK_SO_SNDTIMEO, &forever, sizeof(forever));
\end{ccode}

A timeout of zero means \emph{wait forever}, as with BSD sockets. Sockets created with
\code{zsock_socket()} are not affected: the driver sets their timeouts when it creates them. The
simulator is software, with bugs of its own.

\subsubsection{TCP and UDP}
The \file{lectures/L06/examples/sockets} example is a TCP server: \code{listen()}, \code{accept()} a
connection, then \code{recv()} and \code{send()} a stream of bytes, with no message boundaries,
until the client closes it. UDP, which Exercise Set~\ref{c6:set:5} uses, has no connections: a
server \code{bind()}s a port and answers each datagram, whole, to the address \code{recvfrom()} says
it came from. Nothing is retransmitted, so a client that waits for an answer should wait with a
timeout (\code{ZSOCK_SO_RCVTIMEO}), not forever.

A program using offloaded sockets can be tested from the host with netcat:

\begin{shell}
nc 127.0.0.1 4242       # TCP, the sockets example.
nc -u 127.0.0.1 4300    # UDP, Exercise Set 5.
\end{shell}

\noindent On \code{native_sim}, networking also prints \code{WARNING: Using a test - not safe -
entropy source} at boot: the network stack wants random numbers, and \code{native_sim}’s entropy
device is a seeded generator, deterministic by design.

\subsection{Summary}
\begin{itemize}
  \item The shell is a command line inside the firmware, running in a thread of its own over a backend. On
    \code{native_sim}, \code{CONFIG_UART_NATIVE_PTY_0_ON_STDINOUT} puts it on the terminal.
  \item A command is a handler and a static entry. The mandatory word count includes the command’s name, and
    the shell checks the count; the handler checks what the words say, with \code{strtol()}, not
    \code{atoi()}.
  \item Handlers print with \code{shell_print()} and \code{shell_error()}, and return \code{0} or a negative
    \code{errno}, which \code{retval} shows.
  \item Settings are named key-value pairs, owned by subtree handlers, stored by a backend such as ZMS in
    the flash partition the devicetree names.
  \item \code{settings_load()} calls each handler’s \code{h_set} with the stored length and a function to
    read the value; nothing is loaded until the application asks.
  \item \code{native_sim}’s flash is the file \code{flash.bin} in the working directory, controlled by
    \code{-flash}, \code{-flash_erase}, \code{-flash_rm} and \code{-flash_in_ram}.
  \item Offloaded sockets put the host’s network under Zephyr’s socket API, without setup and without
    Zephyr’s IP stack. Zephyr 4.4 prefixes the socket API’s names with \code{zsock_}, \code{net_} and
    \code{NET_}.
  \item Blocking socket calls block one thread, as long as the socket’s timeouts are set: set them
    explicitly on accepted sockets in Zephyr 4.4.2.
\end{itemize}

% Section 6.2, from lectures/L06/appendix/b_debugging.md.

\section{Debugging: Fatal Errors, Sanitizers, GDB and Coverage}
\label{c6:sec:debuggingfatalerro}\label{c6:app:b}
This section covers what happens when a Zephyr program goes wrong, and the tools that find out why.
Some of them are Zephyr’s own and work the same on a board: fatal errors, the handler that decides
what one costs, and assertions. The others are the host’s, and are much of the reason to keep code
runnable on \code{native_sim}: AddressSanitizer, UndefinedBehaviorSanitizer, GDB and gcov. By the
end you should be able to explain every line the \file{lectures/L06/examples/debugging} example
prints, in each of its configurations.

\subsection{Fatal errors}

\subsubsection{What makes an error fatal}
A \textbf{fatal error} is one the running code cannot continue from. The kernel names the reason
with a number:

\begin{booktable}{@{}llL@{}}
  \thead{Reason} & \thead{Value} & \thead{Raised by} \\
  \midrule
  \code{K_ERR_CPU_EXCEPTION} & 0 & The CPU: a bus fault, an illegal instruction, a division by zero,
    on a board. \\
  \code{K_ERR_SPURIOUS_IRQ} & 1 & An interrupt that has no handler. \\
  \code{K_ERR_STACK_CHK_FAIL} & 2 & A stack overflow check. \\
  \code{K_ERR_KERNEL_OOPS} & 3 & \code{k_oops()}: \emph{this thread} cannot go on. \\
  \code{K_ERR_KERNEL_PANIC} & 4 & \code{k_panic()}, and every failed assertion: \emph{the system}
    cannot go on. \\
\end{booktable}

On \code{native_sim} there is no CPU exception to raise, since the code runs natively; a genuine
segmentation fault kills the host process like any other Linux program’s. Oopses, panics and
assertions work exactly as on a board.

\textbf{\code{k_oops()} and \code{k_panic()} are for code that has found something impossible}, not
for ordinary errors. A sensor that does not answer is an error to return; a reading below absolute
zero, from a sensor the code has already validated, is a broken assumption, and an oops is how a
thread gives up on it.

\subsubsection{What the kernel does}
Every fatal error goes through the same path. The kernel logs what happened, and then calls
\textbf{\code{k_sys_fatal_error_handler()}} with the reason:

\begin{console}
[00:00:00.000,000] <err> app: Impossible reading: -300 C
@ WEST_TOPDIR/lectures/L06/exercises/exercise3/src/main.c:29
[00:00:00.000,000] <err> os: >>> ZEPHYR FATAL ERROR 3: Kernel oops on CPU 0
[00:00:00.000,000] <err> os: Current thread: 0x40f6c0 (sampler)
\end{console}

\begin{itemize}
  \item \textbf{\code{@ file:line}} is printed on architectures that have no CPU exception to raise for an
    oops, \code{native_sim} among them. On a board, the exception’s registers say where instead.
  \item \textbf{\code{ZEPHYR FATAL ERROR 3: Kernel oops}} gives the reason’s number and name. Twister looks
    for this exact string to tell a crashed test from a failed one.
  \item \textbf{\code{Current thread}} names the thread, if \code{CONFIG_THREAD_NAME=y}; otherwise only its
    address.
\end{itemize}

All of it is logged, so \textbf{without \code{CONFIG_LOG=y} a fatal error prints nothing} but what
the handler prints itself.

\subsubsection{The handler}
The kernel’s own \code{k_sys_fatal_error_handler()} is a weak function that halts the system,
whatever the reason:

\begin{console}
[00:00:00.000,000] <err> os: Halting system
Exiting due to fatal error
\end{console}

\noindent On a board, halting means an endless loop with interrupts off, until the watchdog or a
person resets the device. On \code{native_sim}, it means that the process exits with \textbf{status
1}, after printing \code{Exiting due to fatal error} on stderr.

An application that wants a different policy defines the function itself, which replaces the weak
one:

\begin{ccode}
#include <zephyr/fatal.h>
#include <zephyr/logging/log_ctrl.h>

void k_sys_fatal_error_handler(unsigned int reason, const struct arch_esf *esf)
{
	ARG_UNUSED(esf);

	LOG_PANIC();
	LOG_ERR("Fatal error %u in thread '%s'", reason, k_thread_name_get(k_current_get()));

	if (reason == K_ERR_KERNEL_OOPS) {
		return;
	}
	k_fatal_halt(reason);
}
\end{ccode}

\noindent \textbf{What the handler decides is what the error costs:}

\begin{itemize}
  \item \textbf{If it returns}, the kernel aborts only the thread that caused the error, and every other
    thread keeps running. A thread waiting in \code{k_thread_join()} for the aborted one returns,
    because an aborted thread counts as finished.
  \item \textbf{If it calls \code{k_fatal_halt()}}, the system halts, as with the default handler.
\end{itemize}

\textbf{Returning from a panic is not allowed.} With \code{CONFIG_ASSERT=y}, the kernel checks for
it, and the check is itself an assertion, so it panics again, and calls the same handler, which
returns again. On \code{native_sim} that recursion ends when the host process runs out of stack,
with a segmentation fault and exit status 139:

\begin{console}
ASSERTION FAIL [reason != K_ERR_KERNEL_PANIC] @ WEST_TOPDIR/deps/zephyr/kernel/fatal.c:135
	Attempted to recover from a kernel panic condition
\end{console}

\noindent \code{LOG_PANIC()} switches logging to synchronous. On a board, logging is deferred by
default (Chapter~\ref{c1:ch}): the messages wait for a thread of their own, which may never run
again after a fatal error. On \code{native_sim} logging is immediate already, so the call changes
nothing you can see there; it belongs in the handler anyway, because the handler is the same code on
a board.

A real product’s handler usually does more: it saves the reason and the thread’s name somewhere that
survives a reset, such as a retained RAM area or a setting, and resets the device, so that the next
boot can report what happened. The \file{lectures/L06/examples/debugging} example’s handler is
\file{lectures/L06/examples/debugging/src/fatal.c}.

\subsection{Assertions}
An \textbf{assertion} states something the code assumes, and checks it at run time:

\begin{ccode}
__ASSERT((percent >= 0) && (percent <= 100), "duty cycle %d %% is out of range", percent);
\end{ccode}

\noindent \code{__ASSERT_NO_MSG(condition)} is the same without a message. A failed assertion prints
the condition, where it is, and the message, and then panics:

\begin{console}
ASSERTION FAIL [(percent >= 0) && (percent <= 100)] @ WEST_TOPDIR/lectures/L06/exercises/exercise3/src/main.c:40
	duty cycle 150 % is out of range
@ WEST_TOPDIR/deps/zephyr/lib/os/assert.c:43
[00:00:00.000,000] <err> os: >>> ZEPHYR FATAL ERROR 4: Kernel panic on CPU 0
\end{console}

\noindent \textbf{Assertions are compiled out unless \code{CONFIG_ASSERT=y}}, in Zephyr’s code and
in yours. Without it the condition is not even evaluated, and the program above carries on with a
duty cycle of 150 \%. With it, CMake warns that \code{__ASSERT() statements are globally ENABLED},
because every check costs flash and time. That makes assertions a development tool: tests and debug
builds enable them, and most products ship without.

Two consequences follow:

\begin{itemize}
  \item \textbf{An assertion must not have side effects.}
    \code{__ASSERT(k_sem_take(&sem, K_NO_WAIT) == 0, "...")} takes the semaphore only in builds with
assertions.
  \item \textbf{An assertion is not error handling.} Input from a user, a sensor or the network can be wrong
    in a correct program, and has to be checked with an \code{if} in every build.
\end{itemize}

For conditions the compiler can decide, \code{BUILD_ASSERT(condition, "message")} checks at compile
time and costs nothing at run time.

\subsection{Sanitizers}
The \textbf{sanitizers} are compiler instrumentation. The compiler adds a check to every memory
access or arithmetic operation it can, and a runtime library reports every check that fails, with
the source line, the variable, and a stack trace. They find bugs that on a board corrupt memory
silently and crash, if at all, somewhere else, much later.

\subsubsection{AddressSanitizer}
\begin{shell}
west build -p -b native_sim/native/64 lectures/L06/examples/debugging \
    -- -DEXTRA_CONF_FILE=asan.conf
\end{shell}

\code{CONFIG_ASAN=y} builds with \code{-fsanitize=address}. AddressSanitizer surrounds every stack
array, heap block and global with poisoned \emph{redzones}, and checks every read and write against
them. It finds out-of-bounds reads and writes, use after free, and double frees. Exercise
Set~\ref{c6:set:4}’s battery monitor reads one element past an array on the stack:

\begin{console}
==1354852==ERROR: AddressSanitizer: stack-buffer-overflow on address 0x78a35ddf4070 ...
READ of size 4 at 0x78a35ddf4070 thread T2
    #0 0x402a6d in average .../exercise4/src/main.c:18

Address 0x78a35ddf4070 is located in stack of thread T2 at offset 48 in frame
    #0 0x402aa0 in main .../exercise4/src/main.c:33

  This frame has 1 object(s):
    [32, 48) 'window' (line 34) <== Memory access at offset 48 overflows this variable
\end{console}

\noindent Read it from the top: what kind of access (\code{READ of size 4}), where in the code (line
18, in \code{average()}), and whose memory it overflowed (\code{window}, declared on line 34).
\code{T2} is a host thread: every Zephyr thread on \code{native_sim} is one. The report ends with
\code{SUMMARY: AddressSanitizer: stack-buffer-overflow ... in average}, and the process exits with
\textbf{status 1} at the first error. (\code{CONFIG_ASAN_RECOVER}, on by default, compiles the
checks so that they \emph{could} continue, but the runtime only does so when \code{ASAN_OPTIONS}
contains \code{halt_on_error=0}.)

\textbf{AddressSanitizer exists only on \code{native_sim}}: \code{CONFIG_ASAN} depends on
\code{ARCH_POSIX}, because the checks call into the host’s \code{libasan}. The Kconfig help also
notes that leak detection at exit is off for 64-bit boards built with GCC, so
\code{native_sim/native/64} reports no leaks when the program ends.

\subsubsection{UndefinedBehaviorSanitizer}
\begin{shell}
west build -p -b native_sim/native/64 lectures/L06/examples/debugging \
    -- -DEXTRA_CONF_FILE=ubsan.conf
\end{shell}

\code{CONFIG_UBSAN=y} builds with \code{-fsanitize=undefined}, which checks for undefined behaviour
in arithmetic and pointers: signed integer overflow, shifts by too much, division by zero,
misaligned pointers, and out-of-range enum and bool values. Exercise Set~\ref{c6:set:4}’s power
calculation overflows an \code{int32_t}:

\begin{console}
.../exercise4/src/main.c:26:32: runtime error: signed integer overflow: 2250000 * 1000 cannot be represented in type 'int'
5000 mV, average 3725 mV, -2044967296 uW
\end{console}

\noindent Two things differ from AddressSanitizer:

\begin{itemize}
  \item \textbf{The program continues}, with whatever the undefined operation produced, here a wrapped
    negative number. One run can report several problems.
  \item \textbf{The report goes to stderr}, not to the console’s stdout, so a tool that only reads stdout
    never sees it. Twister’s console harness is one such tool.
\end{itemize}

UndefinedBehaviorSanitizer is not tied to \code{native_sim} as tightly. Its reports need a runtime
library, the host’s or picolibc’s; on other targets \code{CONFIG_UBSAN_TRAP} turns undefined
behaviour into a CPU exception, without saying what or where.

\textbf{Both can be enabled at once} (\code{-DCONFIG_ASAN=y -DCONFIG_UBSAN=y}), and a program that
is correct prints exactly the same with them as without. That is how Exercise Set~\ref{c6:set:4} is
tested.

\subsubsection{What a program without them does}
The same battery monitor, built without sanitizers, printed this on the machine these notes were
written on:

\begin{console}
3300 mV, average 249333220 mV, 1485000000 uW
5000 mV, average 249333645 mV, -2044967296 uW
\end{console}

\noindent The out-of-bounds read picked up whatever was next to the array, and the overflow wrapped.
Neither crashed. Another build of the same code averaged \code{261933924 mV} instead, and on another
compiler, or with another optimization level, the average could as easily have looked plausible.
That is the case for the sanitizers: undefined behaviour does not announce itself.

\subsection{GDB}
A \code{native_sim} program is an ordinary Linux executable with debug information, so GDB debugs it
directly, with no probe and no GDB server:

\begin{shell}
sudo apt install gdb
gdb --args ./build/zephyr/zephyr.exe -no-rt
\end{shell}

\begin{aside}
\textbf{Not verified here.} GDB was not installed on the machine these notes were written on, so the
commands in this section were not run while writing them. They are GDB’s standard commands, and the
Kconfig options they rely on were checked in Zephyr’s source.
\end{aside}

The commands that matter most:

\begin{booktable}{@{}LL@{}}
  \thead{Command} & \thead{Effect} \\
  \midrule
  \code{break average}, \code{break main.c:18} & Stop at a function, or at a line. \\
  \code{break k_sys_fatal_error_handler} & Stop at every fatal error, before the handler decides. \\
  \code{run}, \code{continue} & Start the program, or carry on after a stop. \\
  \code{backtrace} & The calls that led to where the program stopped. \\
  \code{print window[i]}, \code{info locals} & Values of variables. \\
  \code{info threads}, \code{thread 2} & The host threads, and switching between them. \\
\end{booktable}

Two Zephyr options make it easier:

\begin{itemize}
  \item \textbf{\code{CONFIG_DEBUG_OPTIMIZATIONS=y}} compiles with \code{-Og}, which the
    \file{lectures/L06/examples/debugging} example sets, so that variables are not optimized away and
    lines run in the order they are written.
  \item \textbf{\code{CONFIG_ARCH_POSIX_TRAP_ON_FATAL=y}} raises \code{SIGTRAP} in the halt path, before the
    process exits, so a debugger stops at a fatal error without a breakpoint.
\end{itemize}

\code{native_sim} creates a host thread for every Zephyr thread when the Zephyr thread is created,
next to the simulator’s own. A thread that is stalled in a busy loop (Chapter~\ref{c3:ch}) is easy
to find this way: interrupt the program with \code{Ctrl+C} in GDB, and look at \code{info threads}
and \code{backtrace}.

\subsection{Coverage}
\textbf{Coverage} measures which lines and branches of the code a run executed. On \code{native_sim}
it is GCC’s own gcov:

\begin{shell}
west build -p -b native_sim/native/64 lectures/L06/exercises/exercise4 \
    -- -DCONFIG_COVERAGE=y
./build/zephyr/zephyr.exe -no-rt -stop_at=1
\end{shell}

\noindent With \code{CONFIG_COVERAGE=y} the program counts every line it runs, and \textbf{writes
the counts when it exits}: \code{-stop_at} is an exit, and so is a fatal error’s. They go into a
\code{.gcda} file next to each object file, such as \code{build/CMakeFiles/app.dir/src/main.c.gcda},
which \code{gcov} reads:

\begin{shell}
cd build
gcov -o CMakeFiles/app.dir/src/main.c.gcda \
    ../lectures/L06/exercises/exercise4/src/main.c
\end{shell}

\begin{console}
Lines executed:100.00% of 14
Creating 'main.c.gcov'
\end{console}

\noindent (The number of lines depends on how your \code{exercise4} is written.)

\code{main.c.gcov} is the source with a count beside every line:

\begin{console}
       24:   18:	for (size_t i = 0; i < n; i++) {
       18:   19:		sum += window[i];
\end{console}

\noindent The loop body ran 18 times over the six calls, \code{1 + 2 + 3 + 4 + 4 + 4}, and its
condition was tested 24 times, once more per call, when it failed.

Twister collects coverage over a whole test run, and turns it into an HTML report:

\begin{shell}
sudo apt install lcov
west twister -T lectures/L06/examples/debugging -p native_sim/native/64 \
    --coverage --coverage-tool lcov
\end{shell}

\noindent The report is \code{twister-out/coverage/index.html}. Twister’s default tool is
\code{gcovr} instead, which is installed with \code{pip install gcovr}.

On a board, coverage is possible too (\code{CONFIG_COVERAGE_GCOV}), but the counts have to be dumped
over a serial line and converted on the host, and the instrumentation needs RAM that small devices
do not have. On \code{native_sim} it is one option.

\subsection{Summary}
\begin{itemize}
  \item A fatal error has a reason: an oops means a thread cannot go on, a panic that the system cannot.
    Assertions panic.
  \item The kernel logs every fatal error, and calls \code{k_sys_fatal_error_handler()}. The default one
    halts; on \code{native_sim} that is exit status 1.
  \item A handler of your own decides what an error costs. Returning aborts only the faulting thread; a
    panic must halt, with \code{k_fatal_halt()}.
  \item \code{__ASSERT()} is compiled out unless \code{CONFIG_ASSERT=y}. Assertions must have no side
    effects, and are not error handling.
  \item AddressSanitizer finds out-of-bounds access and use after free, stops at the first error, and exists
    only on \code{native_sim}.
  \item UndefinedBehaviorSanitizer finds undefined arithmetic, reports on stderr, and continues.
  \item A \code{native_sim} program is debugged with GDB like any Linux program, one host thread per Zephyr
    thread.
  \item \code{CONFIG_COVERAGE=y} writes gcov data when the program exits; Twister’s \code{--coverage} turns
    a test run’s into a report.
\end{itemize}

% Section 6.3, from lectures/L06/README.md: the objectives, the evaluation
% questions, and the next lecture.

\section{Review}
\label{c6:sec:review}

\needspace{6\baselineskip}
\subsection*{What you should be able to do now}
After this chapter, you should:

\begin{itemize}
  \item Be able to add shell commands to an application, and use the shell to inspect a running system.
  \item Be able to store and load settings, and understand where they are stored on \code{native_sim}.
  \item Be able to write TCP and UDP servers with Zephyr’s socket API, and reach them from the host.
  \item Understand how Zephyr handles a fatal error, and be able to decide what happens after one.
  \item Be able to find memory errors and undefined behaviour with the sanitizers.
  \item Be able to debug a \code{native_sim} program with GDB, and measure test coverage.
\end{itemize}

\needspace{6\baselineskip}
\subsection*{Questions to test yourself}
You should be able to explain:

\begin{enumerate}
  \item What does a shell command handler receive, and how does it report an error?
  \item Why can \code{log enable dbg} not show debug messages from a module compiled at level \code{INF}?
  \item Where does the settings subsystem store its data on \code{native_sim}, and what makes a run a first
    boot?
  \item What does the offloaded sockets driver do, and why does it need no network setup or root?
  \item What is the difference between \code{k_oops()} and \code{k_panic()}, and what does
    \code{k_sys_fatal_error_handler()} decide?
  \item Why does AddressSanitizer work on \code{native_sim}, and why is it not available when the same code
    is built for a board?
  \item Why does a UBSan report not stop the program, while an ASan report does?
  \item What would change, and what would not, if the driver from Chapter~\ref{c5:ch} and these subsystems
    moved to a real board?
\end{enumerate}

\needspace{6\baselineskip}
\subsection*{Where to go from here}
This is the last chapter: you have built, configured, described, scheduled, driven, tested and
debugged a Zephyr application, and every one of those programs builds for a real board as it stands.
What is left is the board itself.

\textbf{Try it on hardware.} Any board Zephyr supports will do; two cheap ones its tree already
knows are an STM32 Nucleo (\code{-b nucleo_f401re}) and an ESP32 DevKitC
(\code{-b esp32_devkitc/esp32/procpu}). Three things change in the workspace. You need a
cross-compiler, which \code{west sdk install} fetches, where the host's own GCC has stood in until
now. You need the vendor's HAL: \code{native_sim} needed no Zephyr modules, so the course's
\file{west.yml} imports none, and a real SoC needs its own, such as \code{hal_stm32} or
\code{hal_espressif}. And you flash rather than run: \code{west flash} and \code{west debug}
replace \code{west build -t run}, over the Nucleo's on-board ST-LINK or the DevKitC's USB port.

\textbf{What changes on a board} is the part the simulator was never pretending to be: pins and
clocks have to be configured, the console is a real UART, logging is deferred rather than immediate,
interrupts arrive while your code is running, stack overflows can be caught in hardware, and timing
is real, so a busy loop merely wastes power instead of stopping the world.

\textbf{What does not change} is nearly everything this book is about: the kernel API, devicetree
and Kconfig, the driver model, and your tests, which Twister runs on the board itself with
\code{--device-testing}. That is the point of having written them against an emulator.

% Section 6.4, from lectures/L06/appendix/c_exercises.md.

\section{Exercises}
\label{c6:sec:exercises}\label{c6:app:c}
These exercises reinforce the concepts from \secref{c6:app:a} and \secref{c6:app:b}.

\begin{aside}
\textbf{How to check your work.} Write each exercise set as its own application under
\file{lectures/L06/exercises}, with the file names given there, and run \code{make test} from the
repository root. A set’s tests switch on as soon as its files exist, and a set you have not started
is reported as \code{SKIP}. Each application grows through its set, and the tests check where it
ends up. Shell commands are fed to your program on its standard input, and its output is checked for
the lines each exercise shows, in order; the prompt, the echoed commands and anything else in
between do not matter.

The solutions are published in the companion repository, under
\file{lectures/L06/appendix/solutions}, a few days after the lecture. Try each exercise yourself
before you read them.
\end{aside}

Remember \code{source env.sh} in every new terminal. Build and run an exercise from the repository
root:

\begin{shell}
west build -p -b native_sim/native/64 lectures/L06/exercises/exercise1
west build -t run
\end{shell}

% ---------------------------------------------------------------------
\exerciseset{Shell commands of your own}
\label{c6:set:1}
Create an application named \code{exercise1} to work in for this exercise set.

\exercise{The setpoint}{Code}
\label{c6:ex:1.1}
In this exercise you will give a thermostat a shell command that reads and changes its setpoint.

\task{Tasks}
\begin{itemize}
  \item In \code{prj.conf}, enable the shell, and put the shell’s UART on the terminal the program is
    started from.
  \item In \code{src/main.c}, keep the setpoint in a \code{static int setpoint_c}, initially 21.
  \item Create a root command \code{setpoint} with the help text \code{The thermostat's setpoint.}, no
    handler of its own, and two subcommands:
\begin{itemize}
  \item \code{get}, which takes no arguments, and prints \code{Setpoint: <t> C}.
  \item \code{set <celsius>}, which takes exactly one argument, sets the setpoint, and prints \code{Setpoint
    set to <t> C}.
\end{itemize}
  \item Let the shell check the argument counts: declare them with \code{SHELL_CMD_ARG()}.
  \item \code{main()} prints \code{Thermostat ready} and returns.
\end{itemize}

Pipe a session into the program:

\begin{shell}
printf 'setpoint get\nsetpoint set 25\nsetpoint get\n' \
    | ./build/zephyr/zephyr.exe -stop_at=1
\end{shell}

\noindent The output must contain, in this order:

\begin{console}
Thermostat ready
Setpoint: 21 C
Setpoint set to 25 C
Setpoint: 25 C
\end{console}

\task{Reflection}
\begin{itemize}
  \item The program’s output also contains \code{uart:~$ setpoint get}. Where does that line come from?
  \item Run \code{setpoint} alone, and then \code{retval}. What does \code{retval} print, and why that
    value?
  \item \code{main()} returns at once. Why does the shell keep answering commands?
\end{itemize}

\exercise{Wrong input}{Code}
\label{c6:ex:1.2}

\task{Tasks}
Make \code{setpoint set} reject what is not a valid setpoint:

\begin{itemize}
  \item A word that is not a whole decimal number, such as \code{warm} or \code{25x}, prints the error
    \code{not a temperature: '<word>'} with \code{shell_error()} and returns \code{-EINVAL}.
  \item A number outside 5 to 30 prints the error \code{setpoint must be between 5 and 30 C, not <t>} and
    returns \code{-ERANGE}. The setpoint does not change.
\end{itemize}

For the session

\begin{console}
setpoint get
setpoint set 25
setpoint get
setpoint set 31
retval
setpoint set warm
retval
setpoint set 25x
setpoint set
\end{console}

the output must contain, in this order:

\begin{console}
Setpoint: 21 C
Setpoint set to 25 C
Setpoint: 25 C
setpoint must be between 5 and 30 C, not 31
-34
not a temperature: 'warm'
-22
not a temperature: '25x'
set: wrong parameter count
\end{console}

\noindent \textbf{Hint:} \code{strtol()} from \code{<stdlib.h>} sets its \code{end} argument to the
first character it did not convert.

\task{Reflection}
\begin{itemize}
  \item Why is \code{atoi()} not good enough here?
  \item Which part of the output does the shell print by itself, without your handler being called?
  \item \code{setpoint set 99999999999999999999} is a number too. What does your handler do with it, and
    why?
\end{itemize}

% ---------------------------------------------------------------------
\exerciseset{A setpoint that survives a restart}
\label{c6:set:2}
Copy \code{exercise1} to a new application named \code{exercise2} to work in for this exercise set.

\exercise{Storing it}{Code}
\label{c6:ex:2.1}
In this exercise you will store the setpoint with the settings subsystem, on the simulated flash.

\task{Tasks}
\begin{itemize}
  \item In \code{prj.conf}, enable the settings subsystem with the ZMS backend, and what ZMS needs: flash
    and the flash map.
  \item Keep the setpoint in an \code{int32_t}, stored under the name \code{thermostat/setpoint}.
  \item Register a settings handler for the subtree \code{thermostat} with
    \code{SETTINGS_STATIC_HANDLER_DEFINE()}. Its \code{h_set} callback takes the stored \code{setpoint},
    and rejects a stored value whose length is not that of an \code{int32_t} with \code{-EINVAL}.
  \item In \code{main()}, initialize the settings subsystem, load the settings, and print \code{Loaded
    setpoint: <t> C} instead of \code{Thermostat ready}.
  \item \code{setpoint set} stores the new value with \code{settings_save_one()} before it changes the
    setpoint in RAM. If storing fails, it prints \code{cannot store the setpoint: <error>} with
    \code{shell_error()} and returns the error.
\end{itemize}

Run the program twice, from the same directory:

\begin{shell}
cd build
printf 'setpoint set 25\n' | ./zephyr/zephyr.exe -stop_at=1
printf 'setpoint get\n' | ./zephyr/zephyr.exe -stop_at=1
\end{shell}

\noindent The first run must print \code{Loaded setpoint: 21 C} and \code{Setpoint set to 25 C}, and
the second \code{Loaded setpoint: 25 C} and \code{Setpoint: 25 C}.

\task{Reflection}
\begin{itemize}
  \item Where is the setpoint stored? Find the file, and find the setting’s name in it with \code{strings}.
  \item Run the second command again from the repository root instead of from \code{build/}. What does it
    load, and why?
  \item Run it with \code{-flash_in_ram}. What does it load, and has the file changed?
  \item In \code{h_set}, why read the value into a local variable and check it, rather than straight into
    the setpoint?
\end{itemize}

\exercise{Forgetting it}{Code}
\label{c6:ex:2.2}

\task{Tasks}
Add a subcommand \code{setpoint reset}, which takes no arguments, deletes the stored setpoint with
\code{settings_delete()}, sets the setpoint back to 21, and prints \code{Setpoint reset to 21 C}.

In five runs from the same directory, where each run is given the commands shown, the output must
contain:

\begin{booktable}{@{}lLL@{}}
  \thead{Run} & \thead{Commands} & \thead{Output} \\
  \midrule
  1 & \code{setpoint set 25} & \code{Loaded setpoint: 21 C}, then \code{Setpoint set to 25 C} \\
  2 & \code{setpoint get} & \code{Loaded setpoint: 25 C}, then \code{Setpoint: 25 C} \\
  3 & none, run with \code{-flash_in_ram} & \code{Loaded setpoint: 21 C} \\
  4 & \code{setpoint reset} & \code{Loaded setpoint: 25 C}, then \code{Setpoint reset to 21 C} \\
  5 & none & \code{Loaded setpoint: 21 C} \\
\end{booktable}

\task{Reflection}
\begin{itemize}
  \item Enable \code{CONFIG_SETTINGS_SHELL} on the command line, and run \code{settings list} and
    \code{settings read thermostat/setpoint} after setting the setpoint to 27. What do the four bytes it
    prints mean?
  \item How would you make every run of the tests start from a factory-new device, without deleting files by
    hand?
\end{itemize}

% ---------------------------------------------------------------------
\exerciseset{Fatal errors}
\label{c6:set:3}
Create an application named \code{exercise3} to work in for this exercise set.

\exercise{An impossible reading}{Code}
\label{c6:ex:3.1}
In this exercise you will make a thread give up with \code{k_oops()}, and see what that costs
without a handler of your own.

\task{Tasks}
\begin{itemize}
  \item In \code{prj.conf}, enable logging and thread names.
  \item In \code{src/main.c}, register the log module \code{app}, and define the readings:
\begin{ccode}
static const int readings[] = {21, 22, -300, 23};
\end{ccode}
  \item Write a thread function \code{sampler()} that goes through the readings. A reading below -273 logs
    the error \code{Impossible reading: <t> C} and calls \code{k_oops()}; any other logs \code{Sample
    <t> C}.
  \item In \code{main()}, create the thread with \code{k_thread_create()}, with a 1024-byte stack from
    \code{K_THREAD_STACK_DEFINE()} and priority 5, name it \code{sampler}, and wait for it with
    \code{k_thread_join()}. Then log \code{main() carries on}, and return.
\end{itemize}

Build, run, and check the exit status:

\begin{shell}
./build/zephyr/zephyr.exe -stop_at=1; echo "exit status $?"
\end{shell}

\task{Reflection}
\begin{itemize}
  \item Which lines does the kernel print, and what does each one tell you?
  \item Does \code{main() carries on} appear? What is the exit status, and why?
  \item Build once more with \code{CONFIG_LOG=n} on the command line. What does a fatal error print now?
\end{itemize}

\exercise{A handler of your own}{Code}
\label{c6:ex:3.2}

\task{Tasks}
\begin{itemize}
  \item Create \code{src/fatal.c}, with a log module named \code{fatal}, and define
    \code{k_sys_fatal_error_handler()} in it. It shall:
\begin{itemize}
  \item Switch logging to synchronous with \code{LOG_PANIC()}.
  \item Log the error \code{Fatal error <reason> in thread '<name>'}.
  \item For an oops, log \code{Aborting the thread} and return.
  \item For anything else, log \code{Halting the system} and call \code{k_fatal_halt()}.
\end{itemize}
  \item Add \code{src/fatal.c} to the application in \code{CMakeLists.txt}.
\end{itemize}

The output must now contain, in this order, and the exit status must be 0:

\begin{console}
<inf> app: Sample 21 C
<inf> app: Sample 22 C
<err> app: Impossible reading: -300 C
>>> ZEPHYR FATAL ERROR 3: Kernel oops
<err> fatal: Fatal error 3 in thread 'sampler'
<err> fatal: Aborting the thread
<inf> app: main() carries on
\end{console}

\task{Reflection}
\begin{itemize}
  \item Why does \code{Sample 23 C} never appear?
  \item Why does \code{k_thread_join()} return, when the thread never returned from its function?
  \item Remove the check for the oops, so that the handler always returns. Nothing changes in this program’s
    output. Why not, and why is it still wrong?
\end{itemize}

\exercise{An assertion}{Code}
\label{c6:ex:3.3}

\task{Tasks}
\begin{itemize}
  \item In \code{src/main.c}, add a function:
\begin{ccode}
static void set_heater_duty(int percent)
{
  __ASSERT((percent >= 0) && (percent <= 100), "duty cycle %d %% is out of range", percent);
  LOG_INF("Heater duty cycle %d %%", percent);
}
\end{ccode}
  \item At the end of \code{main()}, after \code{main() carries on}, call \code{set_heater_duty(150)}.
\end{itemize}

Built as usual, the program must also print \code{<inf> app: Heater duty cycle 150 \%} at the end,
and exit with status 0. Built with \code{-DCONFIG_ASSERT=y}, the output must contain, in this order,
and the exit status must be 1:

\begin{console}
<inf> app: main() carries on
duty cycle 150 % is out of range
>>> ZEPHYR FATAL ERROR 4: Kernel panic
<err> fatal: Fatal error 4 in thread 'main'
<err> fatal: Halting the system
\end{console}

\task{Reflection}
\begin{itemize}
  \item Why does the same program pass a duty cycle of 150 \% in one build and halt in the other?
  \item With the handler from the last reflection of Exercise~\ref{c6:ex:3.2}, the one that always returns,
    build with \code{-DCONFIG_ASSERT=y} and run. What happens, and what is the exit status?
  \item Why would \code{__ASSERT(settings_save_one(...) == 0, "...")} be a bug?
\end{itemize}

% ---------------------------------------------------------------------
\exerciseset{Sanitizers}
\label{c6:set:4}
Create an application named \code{exercise4} to work in for this exercise set.

\exercise{A battery monitor}{Code}
\label{c6:ex:4.1}
In this exercise you will find two bugs with the sanitizers, and fix them.

\task{Tasks}
Create \code{CMakeLists.txt}, a \code{prj.conf} with only a comment, and a \code{src/main.c} holding
exactly this:

\begin{ccode}
/*
 * SPDX-License-Identifier: MIT
 *
 * Exercise Set 4: a battery monitor, with two bugs.
 */

#include <zephyr/kernel.h>

/* How many of the most recent readings the average is taken over. */
#define WINDOW 4

/* The average of the first n readings in a window. */
static __noinline int32_t average(const int32_t *window, size_t n)
{
	int32_t sum = 0;

	for (size_t i = 0; i <= n; i++) {
		sum += window[i];
	}
	return sum / (int32_t)n;
}

/* The power drawn at a voltage and a current, in microwatts. */
static __noinline int32_t power_uw(int32_t millivolts, int32_t milliamps)
{
	return millivolts * milliamps * 1000;
}

static const int32_t readings_mv[] = {3300, 3310, 3290, 3300, 5000, 5020};
static const int32_t load_ma = 450;

int main(void)
{
	int32_t window[WINDOW] = {0};

	for (size_t i = 0; i < ARRAY_SIZE(readings_mv); i++) {
		window[i % WINDOW] = readings_mv[i];

		printk("%d mV, average %d mV, %d uW\n", readings_mv[i],
		       average(window, MIN(i + 1, WINDOW)), power_uw(readings_mv[i], load_ma));
	}
	return 0;
}
\end{ccode}

\noindent Build and run it as it is.

\task{Reflection}
\begin{itemize}
  \item Which lines of the output are wrong? Would you have noticed if you did not know what to expect?
  \item The first three lines are right, although the bug in \code{average()} is there from the first call.
    Why do they come out right?
\end{itemize}

\exercise{AddressSanitizer}{Code}
\label{c6:ex:4.2}

\task{Tasks}
Build with \code{-DCONFIG_ASAN=y}, run, and fix the bug that AddressSanitizer reports, changing as
little as possible.

\task{Reflection}
\begin{itemize}
  \item Which line, which variable, and which kind of access does the report name?
  \item What is the exit status, and why do the lines after the report never appear?
\end{itemize}

\exercise{UndefinedBehaviorSanitizer}{Code}
\label{c6:ex:4.3}

\task{Tasks}
Build with \code{-DCONFIG_UBSAN=y}, run, and fix the bug that UndefinedBehaviorSanitizer reports, so
that \code{power_uw()} returns an \code{int64_t} and computes the power without overflow for any
voltage and current that fit in an \code{int32_t} each.

Built with \code{-DCONFIG_ASAN=y -DCONFIG_UBSAN=y}, and built without either, the program must print
exactly:

\begin{console}
3300 mV, average 3300 mV, 1485000000 uW
3310 mV, average 3305 mV, 1489500000 uW
3290 mV, average 3300 mV, 1480500000 uW
3300 mV, average 3300 mV, 1485000000 uW
5000 mV, average 3725 mV, 2250000000 uW
5020 mV, average 4152 mV, 2259000000 uW
\end{console}

\task{Reflection}
\begin{itemize}
  \item The program did not stop at the report. How is that different from AddressSanitizer?
  \item Would \code{return (int64_t)(millivolts * milliamps * 1000);} fix the bug? Try it with UBSan.
  \item Build the fixed program with \code{-DCONFIG_COVERAGE=y}, run it, and run \code{gcov} on
    \code{main.c} as in \secref{c6:app:b}. How many times did the body of the loop in \code{average()}
    run, and why that number?
\end{itemize}

% ---------------------------------------------------------------------
\exerciseset{Sockets}
\label{c6:set:5}
Create an application named \code{exercise5} to work in for this exercise set.

\exercise{A UDP server}{Code}
\label{c6:ex:5.1}
In this exercise you will write a UDP server that answers every datagram in capitals, and a client
for it in the same program, talking through the host’s network with offloaded sockets.

\task{Tasks}
\begin{itemize}
  \item In \code{prj.conf}, enable networking with IPv4 only, sockets, and \code{native_sim}’s offloaded
    sockets, as in the \file{lectures/L06/examples/sockets} example. No logging is needed.
  \item Define \code{PORT} as 4300.
  \item Write a thread \code{server}, with \code{K_THREAD_DEFINE()}, a 2048-byte stack and priority 7,
    which:
\begin{itemize}
  \item Creates a UDP socket, and binds it to \code{PORT} on every address (\code{NET_INADDR_ANY}).
  \item In an endless loop, receives a datagram of up to 64 bytes with \code{zsock_recvfrom()}, turns it
    into capitals with \code{toupper()}, and sends it back with \code{zsock_sendto()} to the address it
    came from.
\end{itemize}
  \item Use the prefixed names throughout: \code{struct net_sockaddr_in}, \code{NET_AF_INET},
    \code{net_htons()}, and so on.
\end{itemize}

\exercise{The client}{Code}
\label{c6:ex:5.2}

\task{Tasks}
In \code{main()}:

\begin{itemize}
  \item Create a UDP socket, and set its receive timeout to one second with \code{ZSOCK_SO_RCVTIMEO}.
  \item Sleep 100 ms, so that the server, which has a lower priority, has bound its port.
  \item For each of the words \code{hello}, \code{zephyr} and \code{native_sim}, send the word to
    \code{127.0.0.1} at \code{PORT}, receive the reply, and print \code{Sent '<word>', received
    '<reply>'}. If no reply arrives, print \code{No reply to '<word>': <errno>} instead.
  \item Close the socket and print \code{Done}.
\end{itemize}

The output must contain, in this order:

\begin{console}
Sent 'hello', received 'HELLO'
Sent 'zephyr', received 'ZEPHYR'
Sent 'native_sim', received 'NATIVE_SIM'
Done
\end{console}

\task{Reflection}
\begin{itemize}
  \item Run the program without \code{-no-rt} and with \code{-stop_at=10}, and send it a datagram from
    another terminal with \code{printf 'from the host' | nc -u -w1 127.0.0.1 4300}. What comes back, and
    what does that say about where the port is?
  \item Why does the client set a receive timeout, when the example’s TCP client does not?
  \item The output also begins with \code{WARNING: Using a test - not safe - entropy source}. Where does
    that come from?
\end{itemize}

